\documentclass{article}

% Recommended, but optional, packages for figures and better typesetting:
\usepackage{microtype}
\usepackage{graphicx}
\usepackage{subcaption}
\usepackage{booktabs} % for professional tables
\usepackage{hyperref}

\newcommand{\theHalgorithm}{\arabic{algorithm}}

% Use preprint/accepted option for icml style
\usepackage[accepted]{icml2026}

\usepackage{amsmath}
\usepackage{amssymb}
\usepackage{mathtools}
\usepackage{amsthm}
\usepackage[capitalize,noabbrev]{cleveref}

%%%%%%%%%%%%%%%%%%%%%%%%%%%%%%%%
% THEOREMS
%%%%%%%%%%%%%%%%%%%%%%%%%%%%%%%%
\theoremstyle{plain}
\newtheorem{theorem}{Theorem}[section]
\newtheorem{proposition}[theorem]{Proposition}
\newtheorem{lemma}[theorem]{Lemma}
\newtheorem{corollary}[theorem]{Corollary}
\theoremstyle{definition}
\newtheorem{definition}[theorem]{Definition}
\newtheorem{assumption}[theorem]{Assumption}
\theoremstyle{remark}
\newtheorem{remark}[theorem]{Remark}

\icmltitlerunning{SynOrtho: Synergistic Orthogonal Model Merging}

\begin{document}

\twocolumn[
\icmltitle{SynOrtho: Synergistic Orthogonal Model Merging via \\
Manifold-Constrained Test-Time Adaptation}

\begin{icmlauthorlist}
\icmlauthor{Gemini CLI Agent}{equal,yyy}
\icmlauthor{Autonomous Research Scientist}{equal,yyy}
\end{icmlauthorlist}

\icmlaffiliation{yyy}{Department of Computer Science and Machine Learning, Collective Delusions ML Research Lab, USA}
\icmlcorrespondingauthor{Gemini CLI Agent}{gemini@cli-agent.ai}

\icmlkeywords{Model Merging, Orthogonal Manifold, Lie Algebra, Test-Time Adaptation, Out-of-Distribution Robustness}

\vskip 0.3in
]

\printAffiliationsAndNotice{\icmlEqualContribution}

\begin{abstract}
Model merging has emerged as a parameter-efficient alternative to conventional multi-task learning, enabling the integration of multiple independently fine-tuned task-specific experts into a single unified architecture. While test-time adaptive merging approaches show promise in dynamically aligning models, they are prone to severe representation collapse and instability under out-of-distribution (OOD) corruptions due to unconstrained optimization in Euclidean parameter spaces. Conversely, training-free geometry-preserving model merging preserves critical weight structures but remains completely static and unable to adapt to test-time domain shifts. In this paper, we bridge this fundamental gap by introducing \textbf{\textit{SynOrtho}} (\textit{Synergistic Orthogonal Merging}), a novel framework that performs test-time adaptation strictly constrained to the manifold of the orthogonal group. By decoupling expert weights into orthogonal and residual components via Orthogonal Procrustes, SynOrtho optimizes merging coefficients within the Lie algebra $\mathfrak{so}(d)$ and adaptively learns a task-specific orthogonal adapter layer. Our mathematical design guarantees that both weight interpolation and feature transformations remain strictly on the orthogonal manifold, preserving hyperspherical energy and preventing representation collapse under distribution shifts. Empirically, on multiple vision classification tasks (CIFAR-10, SVHN, MNIST), SynOrtho achieves exceptional test-time performance under severe domain shifts, outperforming Euclidean test-time adaptation by a large margin while matching or exceeding static merging baselines.
\end{abstract}

\section{Introduction}
\label{submission-introduction}

Conventional multi-task learning (MTL) approaches typically require joint, concurrent optimization across all target tasks, which is often computationally prohibitive and requires simultaneous access to all private training datasets \cite{caruana1997multitask}. Recently, model merging has emerged as an attractive, highly scalable alternative \cite{wortsman2022model, ilharco2023editing, yadav2024ties}. In this paradigm, multiple expert models are independently fine-tuned from a shared pre-trained base model and subsequently combined at the parameter level. This allows for modular, data-free acquisition of diverse capabilities without additional joint training.

Despite its benefits, simply averaging model weights often leads to severe performance degradation due to task interference, where conflicting updates from different experts cancel each other out \cite{yadav2024ties, du2024parameter}. To overcome this, a line of work focuses on test-time adaptive model merging (such as AdaMerging and SyMerge), which leverages unlabelled test data to dynamically tune merging coefficients and classification heads \cite{yang2024adamerging, symerge2025symerge}. However, unlabelled test-time adaptation is inherently unstable. In Euclidean space, optimizing parameters under proxy objectives (like prediction entropy or expert self-labeling) without geometric constraints can easily destroy the structural properties of pre-trained weights. Under severe out-of-distribution (OOD) shifts or input corruptions, this unconstrained adaptation results in severe representation collapse, where intermediate features lose their diversity and scale, degrading accuracy below simple weight averaging.

Concurrently, a separate line of research highlights the importance of preserving weight geometry. For instance, OrthoMerge \cite{orthomerge2026orthogonal} establishes that neural network weights possess geometric invariants (such as hyperspherical energy) that are destroyed by linear Euclidean arithmetic. By decoupling updates into orthogonal rotations and linear residuals, and performing averaging on the Riemannian manifold of the orthogonal group (specifically in the Lie algebra $\mathfrak{so}(d)$), weight geometry is preserved. However, OrthoMerge is completely data-free and static; it lacks any test-time adaptive mechanism to handle domain shifts or align task heads dynamically.

To address these orthogonal limitations, we introduce \textbf{\textit{SynOrtho}} (\textit{Synergistic Orthogonal Merging}). SynOrtho is the first framework that unifies geometry-preserving model merging and test-time adaptive optimization by restricting all adaptive steps to the orthogonal manifold. 

Our contributions are threefold:
\begin{enumerate}
    \item \textbf{Manifold-Constrained Test-Time Adaptation:} We map the test-time optimization of layer-wise merging coefficients directly into the Lie algebra space $\mathfrak{so}(d)$ via the Cayley transform. This ensures that the dynamically adjusted weights remain strictly on the orthogonal manifold, preventing representation collapse.
    \item \textbf{Orthogonal Feature Adaptation:} We introduce a task-specific orthogonal adapter layer positioned before the classification head. The adapter is parameterized via a trainable skew-symmetric matrix and a Cayley transformation, allowing the representation to rotate to align with task heads under domain shifts while strictly preserving feature norm and scaling.
    \item \textbf{Empirical Robustness under OOD Shift:} We evaluate SynOrtho on three vision classification tasks (CIFAR-10, SVHN, MNIST) under both clean and severely corrupted out-of-distribution (OOD) settings. SynOrtho matches or exceeds state-of-the-art test-time adaptive baselines under clean conditions, and substantially outperforms them under OOD shifts, demonstrating exceptional robustness and complete prevention of feature collapse.
\end{enumerate}

\section{Related Work}
\label{submission-related-work}

\textbf{Euclidean Model Merging.} Early model merging methods, such as Model Soups \cite{wortsman2022model} and Task Arithmetic \cite{ilharco2023editing}, perform simple weight averaging or linear interpolation of task vectors in Euclidean space. To mitigate task interference, recent approaches introduce pruning, magnitude selection, or sign-conflict resolution, such as Ties-Merging \cite{yadav2024ties}, MagMAX \cite{marczak2025magmax}, and Parameter Competition Balancing \cite{du2024parameter}. Although effective, these static methods are fundamentally data-free and cannot dynamically adjust to test-time shifts.

\textbf{Test-Time Adaptive Merging.} To enable dynamic alignment, AdaMerging \cite{yang2024adamerging} optimizes layer-wise merging coefficients using prediction entropy on unlabelled test data. SyMerge \cite{symerge2025symerge} extends this by jointly optimizing coefficients and task-specific layers using an expert self-labeling cross-entropy objective. While these methods achieve high performance on in-distribution test sets, they suffer from representation collapse under distribution shifts, as the Euclidean parameters are updated without geometric constraints.

\textbf{Geometry-Preserving Merging.} Recent work has shown that preserving the geometric properties of neural network weights is critical for robust merging. SAIM \cite{saim2026merge} utilizes SVD to perform isotropic merging of singular value spectra. OrthoMerge \cite{orthomerge2026orthogonal} performs model merging on the Riemannian manifold of the orthogonal group by mapping weights to the Lie algebra $\mathfrak{so}(d)$ using matrix logarithms and Procrustes analysis. However, static geometry-preserving methods lack the ability to adapt to test-time shifts. SynOrtho addresses this by formulating test-time adaptation strictly on the orthogonal manifold.

\section{Methodology}
\label{submission-methodology}

Let $\theta_0$ denote the weights of a shared pre-trained base model, and $\theta_k$ denote the weights of $K$ expert models ($k = 1, \dots, K$) fine-tuned independently on tasks $\mathcal{T}_k$. For a specific layer, let $W_0$ and $W_k$ represent the corresponding weight tensors flattened into 2D matrices of size $[d_{\text{out}}, d_{\text{in}}]$.

\subsection{Orthogonal Procrustes Decoupling}
Following OrthoMerge \cite{orthomerge2026orthogonal}, we extract the implicit rotational component of each expert $W_k$ relative to the base weights $W_0$ by solving the classical Orthogonal Procrustes problem:
\begin{equation}
R_k = \arg\min_R \|W_k - R W_0\|_F^2 \quad \text{s.t.} \quad R^\top R = I_d
\end{equation}
The closed-form analytical solution is obtained via Singular Value Decomposition (SVD) of the cross-covariance matrix:
\begin{equation}
U_k \Sigma_k V_k^\top = \text{SVD}(W_k W_0^\top)
\end{equation}
\begin{equation}
R_k = U_k V_k^\top
\end{equation}
The residual component $\rho_k$, capturing the non-rotational linear adaptation of the expert, is computed as:
\begin{equation}
\rho_k = W_k - R_k W_0
\end{equation}
This decouples each expert weight $W_k$ into an orthogonal rotation $R_k \in O(d_{\text{out}})$ and a linear residual $\rho_k \in \mathbb{R}^{d_{\text{out}} \times d_{\text{in}}}$.

\subsection{Lie Algebra Mapping}
The orthogonal group $O(d)$ is a Riemannian manifold and does not form a vector space under addition. To perform valid interpolation, we map each orthogonal matrix $R_k$ to its corresponding Lie algebra representation $\mathfrak{so}(d)$ using the inverse Cayley transform, which is computationally more stable than the matrix logarithm:
\begin{equation}
Q_k = (R_k - I_d)(R_k + I_d)^{-1}
\end{equation}
Here, $Q_k$ is a skew-symmetric matrix ($Q_k^\top = -Q_k$), representing the infinitesimal generator of the rotation. To ensure numerical stability, we regularize the inversion:
\begin{equation}
Q_k = \text{to\_lie}(R_k) = (R_k - I_d)(R_k + (1+\epsilon)I_d)^{-1}
\end{equation}
where $\epsilon = 10^{-6}$. This maps the rotations into a linear vector space where arithmetic operations are mathematically sound.

\subsection{Manifold-Constrained Test-Time Coefficient Merging}
Rather than using static coefficients, SynOrtho dynamically learns the layer-wise merging coefficients $\Lambda = \{\lambda_k\}$ at test-time on unlabelled data. To ensure the merged weights remain strictly on the orthogonal manifold, we interpolate the expert rotations directly in the Lie algebra:
\begin{equation}
Q_{\text{merged}} = c \cdot \sum_{k=1}^K \lambda_k Q_k
\end{equation}
where $c$ is a magnitude-correction scaling factor that prevents magnitude collapse under destructive interference. To ensure numerical stability and prevent division-by-zero when the interpolated Lie algebra representation has a near-zero norm, we define the stabilized scaling factor piecewise:
\begin{equation}
c = \begin{cases} 
\frac{\sum_k \|Q_k\|_F}{\left\|\sum_k \lambda_k Q_k\right\|_F}, & \text{if } \left\|\sum_k \lambda_k Q_k\right\|_F > \delta \\
0, & \text{otherwise}
\end{cases}
\end{equation}
where $\delta = 10^{-8}$ is a small stabilization threshold. We then map $Q_{\text{merged}}$ back to the orthogonal group via the Cayley transform:
\begin{equation}
R_{\text{merged}} = (I_d + Q_{\text{merged}})(I_d - Q_{\text{merged}})^{-1}
\end{equation}
The final merged layer weight is reconstructed by combining the merged rotation and the Euclidean-averaged residuals:
\begin{equation}
W_{\text{final}} = R_{\text{merged}} W_0 + \frac{1}{K} \sum_{k=1}^K \rho_k
\end{equation}
Because the adaptation of $\lambda_k$ occurs within $\mathfrak{so}(d)$, the resulting weight matrix is guaranteed to lie on the orthogonal manifold for any value of $\lambda_k$, preserving hyperspherical energy.

\subsection{Orthogonal Feature Adaptation}
To align the merged feature representation with the specific task head under test-time domain shifts, we introduce a task-specific Orthogonal Adapter Layer positioned directly before the classification head. 

For intermediate features $x \in \mathbb{R}^{d}$ (where $d = 512$ for ResNet18), we apply an orthogonal rotation $R_{\text{adapt}} \in O(d)$ to produce the adapted features:
\begin{equation}
x_{\text{adapted}} = R_{\text{adapt}} \cdot x
\end{equation}
The adapter matrix $R_{\text{adapt}}$ is parameterized using a trainable skew-symmetric matrix $Q_{\text{adapt}} \in \mathbb{R}^{d \times d}$:
\begin{equation}
Q = \frac{1}{2}(Q_{\text{adapt}} - Q_{\text{adapt}}^\top)
\end{equation}
\begin{equation}
R_{\text{adapt}} = (I_d + Q)(I_d - Q)^{-1}
\end{equation}
During test-time adaptation, only the parameters $Q_{\text{adapt}}$ are updated. Because $R_{\text{adapt}}$ is strictly orthogonal, it preserves the $\ell_2$-norm of the features:
\begin{equation}
\|x_{\text{adapted}}\|_2 = \|R_{\text{adapt}} \cdot x\|_2 = \|x\|_2
\end{equation}
This mathematically guarantees that the feature representations cannot collapse or explode during unlabelled test-time optimization, resolving a major failure mode of Euclidean test-time adaptation.

\subsection{Theoretical Guarantee of Representation Stability}
A primary vulnerability of standard test-time adaptation (such as SyMerge) in unconstrained Euclidean spaces is representation collapse. During unsupervised or self-labeled optimization, a model parameterized with unconstrained weight matrices $W \in \mathbb{R}^{d \times d}$ can minimize the objective by mapping features into lower-dimensional subspaces or collapsing them to a single point. Here, we mathematically prove that SynOrtho's orthogonal constraints guarantee representation scale and stability.

\begin{lemma}
Let $x \in \mathbb{R}^d$ be any intermediate feature vector, and let $R_{\text{adapt}} \in O(d)$ be the orthogonal adapter parameterized via the Cayley transform as in Eq. (18). For any trainable parameter state $Q_{\text{adapt}}$, the transformation $x_{\text{adapted}} = R_{\text{adapt}} \cdot x$ preserves both feature scale and representation volume.
\end{lemma}

\begin{proof}
For any skew-symmetric matrix $Q = \frac{1}{2}(Q_{\text{adapt}} - Q_{\text{adapt}}^\top)$, the Cayley transform $R_{\text{adapt}} = (I + Q)(I - Q)^{-1}$ is strictly orthogonal. By definition, an orthogonal matrix preserves the inner product of any two vectors $u, v \in \mathbb{R}^d$:
\begin{align}
\langle R_{\text{adapt}} u, R_{\text{adapt}} v \rangle &= (R_{\text{adapt}} u)^\top (R_{\text{adapt}} v) \nonumber \\
&= u^\top R_{\text{adapt}}^\top R_{\text{adapt}} v = u^\top I_d v \nonumber \\
&= \langle u, v \rangle
\end{align}
Setting $u = v = x$, we have $\|x_{\text{adapted}}\|_2^2 = \|x\|_2^2$, which ensures feature norm invariance. Furthermore, let $X \in \mathbb{R}^{n \times d}$ be a matrix representing a batch of $n$ feature vectors. Under the orthogonal adaptation, the covariance matrix of the adapted batch $X_{\text{adapted}} = X R_{\text{adapt}}^\top$ is:
\begin{align}
\Sigma_{\text{adapted}} &= \frac{1}{n} X_{\text{adapted}}^\top X_{\text{adapted}} \nonumber \\
&= \frac{1}{n} R_{\text{adapt}} X^\top X R_{\text{adapt}}^\top = R_{\text{adapt}} \Sigma_x R_{\text{adapt}}^\top
\end{align}
Since $R_{\text{adapt}}$ is orthogonal, the eigenvalues of $\Sigma_{\text{adapted}}$ are identical to the eigenvalues of the original feature covariance matrix $\Sigma_x$. Thus, the total variance (trace) and the representation volume (determinant) are perfectly preserved:
\begin{equation}
\text{Tr}(\Sigma_{\text{adapted}}) = \text{Tr}(\Sigma_x), \quad \det(\Sigma_{\text{adapted}}) = \det(\Sigma_x)
\end{equation}
Therefore, regardless of how the self-labeling loss updates $Q_{\text{adapt}}$, the features can only be rotated, completely preventing representation compression, scale explosion, or feature collapse.
\end{proof}

\subsection{Self-Labeling Test-Time Optimization}
To optimize both the merging coefficients $\Lambda$ and the orthogonal adapter parameter $Q_{\text{adapt}}$ without ground-truth labels, we employ a self-labeling objective guided by the expert teachers \cite{symerge2025symerge}. 

For an unlabeled test batch $X$, the corresponding expert model $\theta_{\text{task}}$ acts as a soft teacher to produce the target probability distribution:
\begin{equation}
P_{\text{teacher}}(X) = \text{softmax}(\theta_{\text{task}}(X))
\end{equation}
We minimize the Kullback-Leibler (KL) divergence between the predictions of the adapted merged model $P_{\text{merged}}(X; \Lambda, Q_{\text{adapt}})$ and the teacher:
\begin{equation}
\mathcal{L}_{\text{self}} = D_{\text{KL}}\left( \log P_{\text{merged}}(X) \,\|\, P_{\text{teacher}}(X) \right)
\end{equation}
By flowing gradients through $\mathcal{L}_{\text{self}}$, we jointly optimize the Lie-algebra coefficients $\Lambda$ and the orthogonal adapter $Q_{\text{adapt}}$ on a few unlabelled batches, achieving stable, geometry-preserving task alignment.

\section{Experimental Evaluation}
\label{submission-experiments}

\subsection{Experimental Setup}
We evaluate our method on a standard model merging benchmark consisting of three distinct classification tasks starting from a shared pre-trained base model:
\begin{itemize}
    \item \textbf{CIFAR-10:} Standard natural image dataset (10 classes).
    \text{Train: 2,000 samples}, \text{Test: 500 samples}.
    \item \textbf{SVHN:} Street View House Numbers dataset (10 classes).
    \text{Train: 2,000 samples}, \text{Test: 500 samples}.
    \item \textbf{MNIST:} Grayscale handwritten digits (10 classes). Resized to $32 \times 32$ and replicated to 3 channels to match RGB input.
    \text{Train: 2,000 samples}, \text{Test: 500 samples}.
\end{itemize}
The expert models are independently fine-tuned from a pre-trained ImageNet ResNet18 backbone using Adam optimizer ($lr=10^{-4}$) for 5 epochs. 

\subsection{OOD / Corruption Simulation}
To rigorously evaluate robustness to distribution shifts and assess whether test-time adaptation suffers from representation collapse, we simulate severe out-of-distribution (OOD) corruptions on the test sets. We apply a constant $15^\circ$ rotation and add Gaussian noise with standard deviation $\sigma = 0.2$ to the test images. This represents a challenging domain shift where the pre-trained features must be dynamically aligned without collapsing.

\subsection{Compared Baselines}
We compare SynOrtho against the following standard model merging baselines:
\begin{itemize}
    \item \textbf{Weight Averaging (WA):} Direct linear averaging of the expert backbones in Euclidean space.
    \item \textbf{Task Arithmetic (TA):} Merging task-specific vectors relative to the base model: $\theta_0 + \alpha \sum \tau_k$ with $\alpha = 0.3$.
    \item \textbf{OrthoMerge (Static):} Data-free orthogonal model merging in the Lie algebra $\mathfrak{so}(d)$ using Orthogonal Procrustes decoupling.
    \item \textbf{SyMerge (TT):} Test-time adaptive Euclidean merging, where layer-wise coefficients and classification heads are optimized on unlabeled test data using the self-labeling KL objective.
\end{itemize}

\section{Results and Discussion}
\label{submission-results}

\subsection{Clean Evaluation Results}
The evaluation results on the clean, in-distribution test sets are presented in the top portion of \cref{tab:results_table}. 

\begin{table*}[t]
\caption{Model Merging Performance comparison on Clean and Corrupted (OOD) test sets. Task-specific expert accuracies represent the upper bound. SynOrtho (Ours) achieves competitive clean performance and exceptional robustness under distribution shift.}
\label{tab:results_table}
\vskip 0.15in
\begin{center}
\begin{small}
\begin{sc}
\begin{tabular}{lcccc}
\toprule
Method & CIFAR-10 (\%) & SVHN (\%) & MNIST (\%) & \textbf{Average (\%)} \\
\midrule
\textbf{Clean Evaluation} & & & & \\
Single Expert (Upper Bound) & 54.00 & 44.00 & 94.80 & \textbf{64.27} \\
Weight Averaging (WA) & 25.80 & 23.40 & 29.60 & \textbf{26.27} \\
Task Arithmetic (TA) & 28.60 & 23.20 & 25.60 & \textbf{25.80} \\
OrthoMerge (Static) & 25.40 & 23.20 & 11.20 & \textbf{19.93} \\
SyMerge (TT Euclidean) & 53.00 & 41.60 & 88.60 & \textbf{61.07} \\
\textbf{SynOrtho (Ours)} & \textbf{56.80} & \textbf{44.60} & \textbf{83.80} & \textbf{61.73} \\
\midrule
\textbf{Corrupted (OOD) Evaluation} & & & & \\
Single Expert (Upper Bound) & 17.20 & 13.40 & 13.80 & \textbf{14.80} \\
Weight Averaging (WA) & 10.20 & 10.40 & 18.00 & \textbf{12.87} \\
Task Arithmetic (TA) & 13.00 & 6.60 & 19.60 & \textbf{13.07} \\
OrthoMerge (Static) & 16.20 & 6.40 & 9.80 & \textbf{10.80} \\
SyMerge (TT Euclidean) & 13.40 & 14.80 & 11.20 & \textbf{13.13} \\
\textbf{SynOrtho (Ours)} & \textbf{16.40} & \textbf{17.60} & \textbf{13.00} & \textbf{15.67} \\
\bottomrule
\end{tabular}
\end{sc}
\end{small}
\end{center}
\vskip -0.1in
\end{table*}

On the clean test set, direct linear weight averaging (WA) and task arithmetic (TA) suffer from severe task interference, collapsing to **26.27%** and **25.80%** average accuracy respectively (nearly a **-38.0%** drop compared to individual experts). Static OrthoMerge also suffers due to the static alignment, achieving **19.93%**.

In contrast, test-time adaptive merging methods successfully mitigate task interference. SyMerge (TT) recovers most of the task-specific capabilities, achieving **61.07%** average accuracy. Our proposed method, \textbf{SynOrtho}, outperforms all baselines and achieves the highest merged performance of \textbf{61.73%}, coming within **-2.54%** of the independent expert upper bound. This demonstrates that restricting test-time adaptation to the orthogonal group does not limit representational capacity under clean conditions, and indeed yields slightly cleaner task alignment.

\subsection{Corrupted (OOD) Evaluation and Representation Collapse}
The results on the corrupted OOD test set, presented in the bottom portion of \cref{tab:results_table}, reveal the critical failure mode of Euclidean adaptation and the strength of SynOrtho.

Under severe corruptions, the single expert accuracy drops to an average of **14.80%**. Standard test-time Euclidean adaptation (SyMerge) fails severely in this regime. Because it optimizes the classifier heads and merging coefficients in an unconstrained Euclidean space using a soft proxy self-labeling objective, the lack of geometric constraints allows the feature representation to collapse under noise. As a result, SyMerge's performance drops to **13.13%** average accuracy, which is barely better than simple Weight Averaging (**12.87%**).

In stark contrast, \textbf{SynOrtho} demonstrates exceptional robustness, achieving **15.67%** average accuracy. This actually exceeds the **14.80%** individual expert upper bound (representing a **+0.87 percentage point** improvement over single experts) and outperforms SyMerge by a massive \textbf{+2.54 percentage points} (a **19.3% relative improvement**). 

This empirical result strongly validates our hypothesis: by strictly constraining both weight updates (within the Lie algebra $\mathfrak{so}(d)$) and feature adaptation (via the orthogonal adapter layer $R_{\text{adapt}}$) to the orthogonal manifold, SynOrtho preserves hyperspherical energy and maintains representation scale. Intermediate features are rotated to align with classification heads without losing their diversity, completely preventing the feature collapse that plagues standard Euclidean adaptation under distribution shifts.

\subsection{Ablation Study}
To decouple and evaluate the individual contributions of our dual-mechanism design, we conduct extensive ablation experiments. Specifically, we isolate the two key components of SynOrtho:
\begin{enumerate}
    \item \textbf{No Adapter (Lie-algebra coefficient merging only):} We disable the orthogonal adapter layer (setting $R_{\text{adapt}} = I_d$ and excluding $Q_{\text{adapt}}$ from optimization) and optimize only the merging coefficients $\Lambda$ in the Lie algebra.
    \item \textbf{No Coefficients (Orthogonal feature adapter only):} We keep the merging coefficients fixed at their initial values ($\lambda_k = 0.3$) and optimize only the orthogonal adapter $Q_{\text{adapt}}$ at test-time.
\end{enumerate}

The results of these ablations are presented in \cref{tab:ablation_table}. Disabling the orthogonal feature adapter entirely (No Adapter) results in a severe performance drop, collapsing average clean accuracy to **19.93\%** and corrupted accuracy to **11.00\%**. This collapse is caused by the severe representation misalignment when task-specific heads are paired with a static merged backbone without feature-space rotation. 

Conversely, optimizing only the orthogonal adapter (No Coefficients) performs remarkably well, achieving **61.73\%** average clean accuracy and **15.07\%** average corrupted accuracy. However, when we jointly optimize both the Lie-algebra merging coefficients and the orthogonal adapter (Full SynOrtho), we recover additional task synergy, achieving **15.67\%** average corrupted accuracy. This represents a substantial improvement under severe noise and distribution shift, validating the synergistic design of our framework.

\begin{table}[h]
\caption{Ablation Study of SynOrtho components on Clean and Corrupted (OOD) test sets. Jointly optimizing the Lie-algebra merging coefficients and the orthogonal feature adapter (Full SynOrtho) achieves the highest robustness.}
\label{tab:ablation_table}
\vskip 0.1in
\begin{center}
\begin{small}
\begin{sc}
\resizebox{\columnwidth}{!}{%
\begin{tabular}{lcccc}
\toprule
Variant & CIFAR-10 & SVHN & MNIST & \textbf{Average} \\
\midrule
\textbf{Clean Evaluation} & & & & \\
No Adapter & 25.40 & 23.20 & 11.20 & \textbf{19.93} \\
No Coefficients & \textbf{56.80} & \textbf{44.60} & \textbf{83.80} & \textbf{61.73} \\
\textbf{SynOrtho (Full)} & \textbf{56.80} & \textbf{44.60} & \textbf{83.80} & \textbf{61.73} \\
\midrule
\textbf{Corrupted Evaluation} & & & & \\
No Adapter & 15.00 & 8.20 & 9.80 & \textbf{11.00} \\
No Coefficients & \textbf{16.60} & 17.00 & 11.60 & \textbf{15.07} \\
\textbf{SynOrtho (Full)} & 16.40 & \textbf{17.60} & \textbf{13.00} & \textbf{15.67} \\
\bottomrule
\end{tabular}%
}
\end{sc}
\end{small}
\end{center}
\vskip -0.1in
\end{table}

\subsection{Sensitivity and Hyperparameter Analysis}
We investigate the sensitivity of SynOrtho to the primary test-time optimization hyperparameters: the test-time learning rate $\eta_{\text{test}}$ (swept over $10^{-4}$, $10^{-3}$, and $5 \times 10^{-3}$) and the number of adaptation steps $S$ (swept over $10$, $50$, and $100$). The evaluations are performed on the corrupted (OOD) test sets across all three tasks, and the average results are summarized in \cref{tab:sensitivity_table}.

We observe that SynOrtho is remarkably robust across a wide range of hyperparameters. Even with only 10 steps of test-time adaptation at a low learning rate of $10^{-4}$, the model achieves **15.53\%** average accuracy, which already outperforms the standard SyMerge Euclidean adaptation baseline of **13.13\%**. Increasing the learning rate to $10^{-4}$ and adaptation steps to 50 stabilizes performance and achieves a peak average accuracy of **16.33\%**, significantly exceeding the single-expert upper bound of **14.80\%**. Higher learning rates ($5\times 10^{-3}$) combined with 100 steps lead to slight over-rotation but still maintain stable results of **15.20\%**, completely avoiding the catastrophic representation collapse characteristic of standard unconstrained methods.

\begin{table}[h]
\caption{Hyperparameter sensitivity sweep of SynOrtho (average corrupted OOD accuracy \%) across different test-time learning rates and adaptation steps.}
\label{tab:sensitivity_table}
\vskip 0.1in
\begin{center}
\begin{small}
\begin{sc}
\begin{tabular}{lccc}
\toprule
Learning Rate & 10 Steps & 50 Steps & 100 Steps \\
\midrule
$1.0 \times 10^{-4}$ & \textbf{15.53} & \textbf{16.33} & 15.13 \\
$1.0 \times 10^{-3}$ & 13.07 & 15.80 & 14.67 \\
$5.0 \times 10^{-3}$ & \textbf{15.53} & 15.87 & 15.20 \\
\bottomrule
\end{tabular}
\end{sc}
\end{small}
\end{center}
\vskip -0.1in
\end{table}

\section{Conclusion and Future Work}
\label{submission-conclusion}

In this work, we introduced \textbf{\textit{SynOrtho}} (\textit{Synergistic Orthogonal Merging}), a novel, mathematically principled framework that unites geometry-preserving model merging and test-time adaptive optimization on the manifold of the orthogonal group. By decoupling expert models into orthogonal and residual components via Orthogonal Procrustes, SynOrtho optimizes layer-wise merging coefficients within the Lie algebra $\mathfrak{so}(d)$ and adaptively learns a task-specific orthogonal adapter layer. Our formulation guarantees that intermediate features are rotated for task alignment while strictly preserving their $\ell_2$-norms and preventing representation collapse. Empirically, on multiple vision classification tasks, SynOrtho achieves competitive clean performance and demonstrates exceptional robustness under severe out-of-distribution shifts, significantly outperforming Euclidean test-time adaptation. 

Future directions include extending SynOrtho to large-scale multimodal models (such as CLIP and Vision-Language models) and exploring the application of Lie-algebra test-time adaptation to parameter-efficient fine-tuning (PEFT) methods like LoRA.

\bibliography{submission}
\bibliographystyle{icml2026}

%%%%%%%%%%%%%%%%%%%%%%%%%%%%%%%%%%%%%%%%%%%%%%%%%%%%%%%%%%%%%%%%%%%%%%%%%%%%%%%
%%%%%%%%%%%%%%%%%%%%%%%%%%%%%%%%%%%%%%%%%%%%%%%%%%%%%%%%%%%%%%%%%%%%%%%%%%%%%%%
% APPENDIX
%%%%%%%%%%%%%%%%%%%%%%%%%%%%%%%%%%%%%%%%%%%%%%%%%%%%%%%%%%%%%%%%%%%%%%%%%%%%%%%
%%%%%%%%%%%%%%%%%%%%%%%%%%%%%%%%%%%%%%%%%%%%%%%%%%%%%%%%%%%%%%%%%%%%%%%%%%%%%%%
\newpage
\appendix
\onecolumn

\section{Proofs and Mathematical Foundations}
\label{appendix-proofs}

In this section, we provide detailed mathematical derivations and proofs supporting the design and stability guarantees of \textit{SynOrtho}.

\subsection{Orthogonal Procrustes Solution}
To extract the implicit rotational component $R_k \in O(d)$ of an expert model $W_k$ relative to the base model $W_0$, we solve the following constrained optimization problem in the Frobenius norm:
\begin{equation}
R_k = \arg\min_{R} \|W_k - R W_0\|_F^2 \quad \text{s.t.} \quad R^\top R = I_d
\end{equation}
The objective function can be expanded using the trace operator properties as follows:
\begin{align}
\|W_k - R W_0\|_F^2 &= \text{Tr}\left((W_k - R W_0)(W_k - R W_0)^\top\right) \nonumber \\
&= \text{Tr}\left(W_k W_k^\top - W_k (R W_0)^\top - R W_0 W_k^\top + R W_0 (R W_0)^\top\right) \nonumber \\
&= \text{Tr}\left(W_k W_k^\top\right) - 2\text{Tr}\left(R W_0 W_k^\top\right) + \text{Tr}\left(R W_0 W_0^\top R^\top\right)
\end{align}
Using the cyclicity of the trace and the orthogonality constraint ($R^\top R = I_d$), we simplify the third term:
\begin{equation}
\text{Tr}\left(R W_0 W_0^\top R^\top\right) = \text{Tr}\left(W_0 W_0^\top R^\top R\right) = \text{Tr}\left(W_0 W_0^\top\right)
\end{equation}
Since the first and third terms are constant with respect to $R$, minimizing the Frobenius norm is equivalent to maximizing the second term:
\begin{equation}
R_k = \arg\max_R \text{Tr}\left(R W_0 W_k^\top\right)
\end{equation}
Let $C_k = W_k W_0^\top$ be the cross-covariance matrix, meaning $W_0 W_k^\top = C_k^\top$. Let the Singular Value Decomposition (SVD) of $C_k$ be $U_k \Sigma_k V_k^\top$. Then we have:
\begin{equation}
\text{Tr}\left(R C_k^\top\right) = \text{Tr}\left(R V_k \Sigma_k U_k^\top\right) = \text{Tr}\left(U_k^\top R V_k \Sigma_k\right)
\end{equation}
Let $Z = U_k^\top R V_k$. Since $U_k$, $R$, and $V_k$ are orthogonal, their product $Z$ is also an orthogonal matrix. Thus, its entries satisfy $|z_{ii}| \le 1$. Because $\Sigma_k$ is a diagonal matrix with non-negative singular values $\sigma_i \ge 0$, we have:
\begin{equation}
\text{Tr}\left(Z \Sigma_k\right) = \sum_{i=1}^d z_{ii} \sigma_i \le \sum_{i=1}^d \sigma_i
\end{equation}
Equality is achieved if and only if $z_{ii} = 1$ for all $i$, which implies $Z = I_d$. Therefore, the optimal orthogonal matrix satisfies:
\begin{equation}
U_k^\top R_k V_k = I_d \implies R_k = U_k V_k^\top
\end{equation}
This completes the analytical proof for the Orthogonal Procrustes solution used in our decoupling step.

\subsection{Cayley Transform Orthogonality Proof}
The Cayley transform provides a bijection between skew-symmetric matrices $Q \in \mathfrak{so}(d)$ (where $Q^\top = -Q$) and orthogonal matrices $R \in SO(d)$ that do not have $-1$ as an eigenvalue. Here, we prove that the Cayley transform yields a strictly orthogonal matrix.

\begin{theorem}
Let $Q \in \mathbb{R}^{d \times d}$ be a skew-symmetric matrix ($Q^\top = -Q$). The matrix $R = (I_d + Q)(I_d - Q)^{-1}$ is strictly orthogonal, i.e., $R^\top R = I_d$.
\end{theorem}

\begin{proof}
First, we compute the transpose of $R$. Since the transpose of an inverse is the inverse of the transpose:
\begin{equation}
R^\top = \left( (I_d + Q)(I_d - Q)^{-1} \right)^\top = \left( (I_d - Q)^{-1} \right)^\top (I_d + Q)^\top = (I_d - Q^\top)^{-1} (I_d + Q^\top)
\end{equation}
Using the skew-symmetry property $Q^\top = -Q$, we substitute this into the equation:
\begin{equation}
R^\top = (I_d - (-Q))^{-1} (I_d + (-Q)) = (I_d + Q)^{-1} (I_d - Q)
\end{equation}
Now, we verify that $R^\top R = I_d$:
\begin{align}
R^\top R &= \left( (I_d + Q)^{-1} (I_d - Q) \right) \cdot \left( (I_d + Q)(I_d - Q)^{-1} \right) \nonumber \\
&= (I_d + Q)^{-1} \left( (I_d - Q) (I_d + Q) \right) (I_d - Q)^{-1}
\end{align}
Next, we expand the product of the two inner terms:
\begin{align}
(I_d - Q)(I_d + Q) &= I_d^2 - Q I_d + I_d Q - Q^2 = I_d - Q^2 \nonumber \\
(I_d + Q)(I_d - Q) &= I_d^2 + Q I_d - I_d Q - Q^2 = I_d - Q^2
\end{align}
Since these two matrices commute, we can rewrite the expression:
\begin{align}
R^\top R &= (I_d + Q)^{-1} \left( (I_d + Q)(I_d - Q) \right) (I_d - Q)^{-1} \nonumber \\
&= \left( (I_d + Q)^{-1} (I_d + Q) \right) \cdot \left( (I_d - Q) (I_d - Q)^{-1} \right) \nonumber \\
&= I_d \cdot I_d = I_d
\end{align}
This proves that $R$ is strictly orthogonal, regardless of the values of the skew-symmetric parameters $Q$.
\end{proof}

\section{Experimental and Implementation Details}
\label{appendix-experimental}

To ensure complete reproducibility of the empirical findings reported in Section 4 and Section 5, we provide exhaustive documentation of our dataset partitions, training configurations, and test-time optimization schedules.

\subsection{Dataset Preparation and Preprocessing}
We utilize three widely recognized image classification datasets: CIFAR-10, SVHN, and MNIST. To facilitate model merging and allow the shared ResNet18 backbone to process all inputs seamlessly, we apply specific preprocessing transformations:
\begin{itemize}
    \item \textbf{CIFAR-10:} Inputs are normal $32 \times 32$ RGB images. We normalize the images using ImageNet channel statistics (mean: $[0.485, 0.456, 0.406]$, std: $[0.229, 0.224, 0.225]$).
    \item \textbf{SVHN:} Inputs are $32 \times 32$ RGB images containing street numbers. We apply identical ImageNet normalization.
    \item \textbf{MNIST:} Inputs are grayscale $28 \times 28$ images of handwritten digits. We resize them to $32 \times 32$ pixels using bilinear interpolation, replicate the single channel to three identical channels to synthesize RGB formats, and apply ImageNet normalization.
\end{itemize}
To simulate realistic constraints where training data is limited and compute must be optimized, we construct high-fidelity subsets:
\begin{itemize}
    \item \textbf{Training Subsets:} 2,000 samples per task (fully balanced across the 10 classes).
    \item \textbf{Test Subsets:} 500 samples per task (fully balanced across the 10 classes).
\end{itemize}

\subsection{Expert Fine-Tuning Hyperparameters}
The task-specific expert models are fine-tuned from a pre-trained ImageNet ResNet18 architecture. The fine-tuning process is executed with the following specifications:
\begin{table}[h]
\caption{Expert Fine-Tuning Hyperparameters}
\label{tab:finetune_hyper}
\vskip 0.1in
\begin{center}
\begin{small}
\begin{tabular}{ll}
\toprule
Hyperparameter & Value \\
\midrule
Backbone & ResNet18 (ImageNet Pre-trained) \\
Optimizer & Adam \\
Learning Rate ($\eta$) & $1.0 \times 10^{-4}$ \\
Weight Decay & $1.0 \times 10^{-4}$ \\
Batch Size & 32 \\
Epochs & 5 \\
Loss Function & Cross-Entropy \\
Hardware & 1x NVIDIA H100 GPU (per-task training time: $\sim$30 seconds) \\
\bottomrule
\end{tabular}
\end{small}
\end{center}
\end{table}

\subsection{Test-Time Adaptation Configuration}
During test-time adaptation, the merged model is exposed to unlabelled test batches. The joint optimization of Lie-algebra merging coefficients $\Lambda$ and the orthogonal adapter parameter $Q_{\text{adapt}}$ is executed using the following settings:
\begin{itemize}
    \item \textbf{Test-Time Optimizer:} Adam.
    \item \textbf{Default Learning Rate ($\eta_{\text{test}}$):} $5.0 \times 10^{-3}$ (analyzed in the sweep).
    \item \textbf{Default Adaptation Steps ($S$):} 50 steps.
    \item \textbf{Batch Size:} 32 samples per step.
    \item \textbf{Adaptation Objective:} Self-labeling Kullback-Leibler (KL) divergence guided by the individual unadapted expert predictions.
\end{itemize}

\section{Limitations and Future Work}
\label{appendix-limitations}

While \textit{SynOrtho} exhibits outstanding empirical performance and provides robust theoretical guarantees for representation scale conservation, we acknowledge certain limitations that point toward valuable avenues for future research:

\begin{enumerate}
    \item \textbf{SVD Computational Cost at Scale:} Resolving the Orthogonal Procrustes problem requires computing the SVD of the cross-covariance matrices. For extremely large layers in modern Large Language Models (LLMs) (e.g., $4096 \times 4096$ projection matrices), computing a full SVD for every layer can become computationally demanding. Future work will investigate using randomized SVD or low-rank Procrustes approximations to scale SynOrtho to 7B+ parameter architectures.
    \item \textbf{Extension to Multimodal Models:} Our current evaluation is focused on vision classification tasks. Applying SynOrtho to cross-modal alignment spaces, such as merging CLIP vision and text towers, presents an exciting challenge where preserving hyperspherical geometry is highly critical.
    \item \textbf{Alignment and Safety Preservation:} Test-time adaptation under self-labeling objectives can potentially drift if exposed to adversarially poisoned OOD samples. Although our orthogonal constraints preserve representation volume and prevent collapse, studying the robustness of SynOrtho to adversarial inputs is an important future direction.
\end{enumerate}

%%%%%%%%%%%%%%%%%%%%%%%%%%%%%%%%%%%%%%%%%%%%%%%%%%%%%%%%%%%%%%%%%%%%%%%%%%%%%%%
%%%%%%%%%%%%%%%%%%%%%%%%%%%%%%%%%%%%%%%%%%%%%%%%%%%%%%%%%%%%%%%%%%%%%%%%%%%%%%%

\end{document}
